% ============================================================================
% CAPÍTULO: V762 → V764 — LA AUDITORÍA DE 12 HALLAZGOS Y EL ENDURECIMIENTO
% ============================================================================
\chapter{V762 a V764: La Auditoría de 12 Hallazgos y el Endurecimiento Final}


\textit{``El documento de entrega certificaba únicamente que dlopen y dlsym funcionaron.\\
El Exit Code 0 de V761 era la historia más costosa de POLYDIM:\\
certificaba el trabajo menos relevante.''} --- --- Comentario interno kernel\_cpp\_v764.cpp, línea 6

\section{Contexto: El Día Después de la Certificación V762}

El 19 de septiembre de 2026, a las 18:00 horas (Argentina), POLYDIM V762 alcanzó
Exit Code 0 en 5/5 suites adversariales. Al día siguiente, el 20 de septiembre,
una nueva ronda de auditoría adversarial --- la auditoría más implacable de toda
la historia del proyecto --- identificó \textbf{12 hallazgos críticos} (A1--A12)
que, aunque no comprometían los resultados numéricos reportados, sí constituían
deficiencias en la \textit{completitud} del contrato de correctitud del kernel.

Este capítulo documenta sistemáticamente los 12 hallazgos, sus correcciones,
y las lecciones que cada uno aporta a la teoría del diseño de sistemas numéricos
de alta confiabilidad.

\textbf{[CRÍTICO]} 
\textbf{Nota Importante:} Los benchmarks de V762 (Drift $= 4.44 \times 10^{-16}$,
PMTP MaxBitDiff $= 0$, etc.) son correctos y no son afectados por los hallazgos A1--A12.
Los hallazgos son fallas de \textit{completitud de la API} y \textit{robustez ante entradas
inusuales}, no errores en los resultados certificados.


\section{A1: PMTP Reescrito como Triple Buffer con SEQLock por Ranura}

\subsection{El Problema}

El PMTP de V762 implementaba un doble buffer. El escritor producía en el buffer
inactivo y publicaba atómicamente el puntero al buffer activo. El problema es que
en un sistema con múltiples lectores lentos y un escritor rápido, puede ocurrir:

\begin{enumerate}
\item Escritor: publica en buffer 0 (activo = 0).
\item Lector lento: inicia lectura de buffer 0.
\item Escritor: produce en buffer 1, publica (activo = 1).
\item Escritor: produce nuevamente en buffer 0, publica (activo = 0).
\item Lector lento: \textbf{todavía leyendo buffer 0, que acaba de ser sobreescrito}.
\end{enumerate}

\subsection{La Solución V764}

V764 introduce un SEQLock \textit{por ranura} (por slot) con verificación post-lectura:

\begin{verbatim}
struct PmtpSlot {
    alignas(64) std::atomic<uint64_t> seq; // Impar = escritura en progreso
    // ... datos del tensor ...
};

// ESCRITURA:
// 1. Seleccionar slot libre: not == lector activo
// 2. seq++ (hacer impar: señaliza escritura)
// 3. Escribir datos
// 4. seq++ (hacer par: publicar)

// LECTURA:
// 1. Leer seq1 = slot.seq.load(Acquire)
// 2. Si seq1 impar: slot en escritura, esperar
// 3. Leer datos
// 4. Leer seq2 = slot.seq.load(Acquire)
// 5. Si seq1 != seq2: datos inconsistentes, reintentar

// polydim_pmtp_validate_read: verifica que la secuencia no cambió
int32_t polydim_pmtp_validate_read(PmtpControl* ctrl, uint64_t slot, uint64_t ticket) {
    if (!ctrl) return POLYDIM_ERR_NULL_POINTER;
    // Ticket = valor de seq observado antes de leer
    // Slot actual: verificar que seq sigue igual
    uint64_t current_seq = ctrl->slots[slot].seq.load(std::memory_order_acquire);
    if (current_seq != ticket) return POLYDIM_ERR_SEQLOCK_RACE;
    return POLYDIM_SUCCESS;
}
\end{verbatim}

\begin{theorem}[Correctitud del PMTP Triple Buffer]

El protocolo PMTP V764 con triple buffer y SEQLock por ranura es seguro contra
\textit{reader-writer aliasing}: ningún lector puede observar un buffer durante
su sobreescritura, siempre que existan al menos 3 ranuras.
\end{theorem}

\begin{proof}
Con 3 ranuras y un escritor que nunca sobreescribe la ranura en uso por el lector:
el escritor tiene al menos 1 ranura libre (la que no está siendo leída ni la activa anterior).
El protocolo SEQLock por ranura garantiza que si el escritor inicia una escritura
en la ranura siendo leída (imposible por diseño, pero verificable), el ticket
del lector cambiará y \texttt{pmtp\_validate\_read} retornará \texttt{SEQLOCK\_RACE}.
\end{proof}

\section{A2: Compuerta de Ortonormalidad --- check\_basis()}

\subsection{El Problema}

En V762, el kernel \textbf{calculaba} las normas de $\mathbf{u}$, $\mathbf{v}$ y
$\langle \mathbf{u}, \mathbf{v} \rangle$ en la Pasada 1, pero no usaba esos valores
para validar que la base fuera ortonormal antes de aplicar la rotación.

El resultado: si el usuario pasaba $\mathbf{u}$ y $\mathbf{v}$ no ortonormales,
el kernel devolvía \texttt{SUCCESS} con un resultado incorrecto en lugar de rechazar
la entrada explícitamente.

\subsection{La Solución V764}

\begin{verbatim}
/* A2: la compuerta que V762 calculaba y descartaba. Coste medido: ~1%. */
if (e_uu > tol.basis_ortho || e_vv > tol.basis_ortho || e_uv > tol.basis_ortho)
    return POLYDIM_ERR_BASIS_NOT_ORTHONORMAL;
// Nuevo código de error: -9
// Mensaje: "BASIS_NOT_ORTHONORMAL"
// Solución para el llamante: usar polydim_orthonormalize_pair_f64 primero
\end{verbatim}

\begin{remark}
El hallazgo A2 demuestra un patrón de error común en el diseño de APIs numéricas:
calcular información de validación pero descartarla antes de usarla. El costo de
la validación ya está pagado (los productos punto se calculan en la Pasada 1 de
todas formas). No usarlos para validar es un regalo al bug.
\end{remark}

\section{A3: Validación de theta/tau --- require\_finite()}

\subsection{El Problema}

V762 no validaba que el ángulo $\theta$ fuera finito antes de usarlo en
\texttt{std::sin(theta)} y \texttt{std::cos(theta/2)}.

Con $\theta = \text{NaN}$: \texttt{sin(NaN) = NaN}, \texttt{cos(NaN) = NaN}.
Los coeficientes $\alpha = \text{NaN}$, $\beta = \text{NaN}$ producen
\texttt{y\_out[i] = NaN} para todo $i$. La verificación post-rotación
$\|{y_{\text{out}\|}} = \text{NaN}$, que se comparaba con la tolerancia:
\texttt{NaN > tol = false}. Resultado: \texttt{SUCCESS} con vector de salida NaN.

\subsection{La Solución V764}

\begin{verbatim}
inline bool require_finite(double x) { return std::isfinite(x); }

// Primeras líneas de polydim_rodrigues_geodesic_f64:
if (!require_finite(theta)) return POLYDIM_ERR_INVALID_SCALAR; // -8 (nuevo)
\end{verbatim}

\begin{proposition}[NaN-Transparency de Comparaciones]
Para cualquier $x \in \mathbb{R}$: \texttt{NaN > x} $\equiv$ \texttt{false},
\texttt{NaN < x} $\equiv$ \texttt{false}, \texttt{NaN == x} $\equiv$ \texttt{false}.
En particular, \texttt{NaN > tol} $\equiv$ \texttt{false} para cualquier tolerancia finita,
haciendo que toda verificación de tolerancia con NaN falle silenciosamente.
\end{proposition}

Esto es IEEE-754 §6.2: las comparaciones de orden producen \texttt{false} cuando
cualquiera de los operandos es NaN (excepto $\ne$ que produce \texttt{true}).

\section{A4: Validación del Punto de Entrada Contra la Variedad}

\subsection{El Problema}

V762 no verificaba que $\mathbf{y}$ ya estuviera en $S^{D-1}$ antes de aplicar la rotación.
Si el usuario pasaba un vector con $\|{\mathbf{y}\|} = 1.5$ (no normalizado), el kernel
devolvía un resultado incorrecto con SUCCESS.

\subsection{La Solución V764}

\begin{verbatim}
/* A4: el punto debe estar sobre la variedad. */
if (e_yy > tol.point_norm) return POLYDIM_ERR_POINT_OFF_MANIFOLD; // -10

/* tol.point_norm = 64 * eps = 1.42e-14 (NO escala con D) */
/* Consecuencia: el llamante DEBE normalizar con polydim_project_sphere_f64 */
\end{verbatim}

\section{A5: Tolerancias que NO Escalan con D --- La Decisión de Diseño más Importante}

\subsection{El Error de V762}

V762 (y su documentación) reportó que la tolerancia de Higham escala como:
\begin{equation}
\text{tol}(D) = 2D \cdot \varepsilon_{\text{mach}} + 50\varepsilon_{\text{mach}}
\end{equation}

Esta es la cota \textit{teórica} de Higham para sumatoria \textit{sin} compensación.
El kernel V762 usaba sumación de Neumaier, lo que reduce el error a
$\mathcal{O}(\varepsilon_{\text{mach}})$ \textbf{independiente de $D$}.

El guardián de V762 usaba la cota escalada ($\approx 4.44 \times 10^{-10}$ para $D=10^6$)
cuando el error real medido era $4.44 \times 10^{-16}$ --- es decir, el guardián era
\textbf{$2.1 \times 10^{6}$ veces más laxo que la realidad}.

\subsection{La Corrección V764: El Comentario más Importante del Kernel}

\begin{verbatim}
extern "C" POLYDIM_EXPORT PolydimTolerances POLYDIM_CALL
polydim_default_tolerances(uint64_t D) {
    PolydimTolerances t;
    /* A5 — DECISIÓN DE DISEÑO, no un detalle de implementación.
     *
     * Con sumación compensada la deriva medida es O(eps) e INDEPENDIENTE de D:
     * verificado a 0.00e+00 en D=1e6 vía polydim_project_sphere_f64. Por tanto
     * la cota por defecto NO escala con D. Escalarla fue exactamente el error
     * de V762: su guardián admitía 4.44e-10 en D=1e6, 2.1e6x más laxo que el
     * 2.10e-14 que el documento de entrega declaraba certificado.
     *
     * 64*eps = 1.42e-14 <= 2.10e-14, o sea la cota es MÁS ESTRICTA que la cifra
     * publicada. Consecuencia para el llamante: los vectores deben normalizarse
     * con polydim_project_sphere_f64 (compensado). Una normalización ingenua en
     * D=1e6 deja un error relativo de ~4.3e-14 y será RECHAZADA, con razón.
     */
    t.basis_ortho      = 64.0 * kEps;           /* 1.42e-14, fijo */
    t.point_norm       = 64.0 * kEps;           /* 1.42e-14, fijo */
    t.gram_ortho       = 64.0 * kEps * std::sqrt((double)D); /* Stiefel escala */
    t.pivot_rel        = 8.0  * kEps;
    t.reject_subnormal = 0;
    return t;
}
\end{verbatim}

\begin{theorem}[Independencia de $D$ del Error de Neumaier]

El acumulador de Neumaier satisface:
\begin{equation}
|{\text{NeumaierSum}|(a_1, \ldots, a_D) - \sum_{i=1}^D a_i} \le (2\varepsilon_{\text{mach}} + \mathcal{O}(\varepsilon_{\text{mach}}^2)) \cdot \max_i |a_i|
\end{equation}
\textbf{independiente de $D$}. La cota de Higham $\tilde{u}_D = (2D + 50)\varepsilon_{\text{mach}}$
se aplica a la sumatoria naïve sin compensación y es innecesariamente laxa para Neumaier.
\end{theorem}

La tolerancia correcta es $64\varepsilon_{\text{mach}} = 1.42 \times 10^{-14}$
(64 operaciones de redondeo en el peor caso del bucle compensado), no $2D\varepsilon_{\text{mach}}$.

\section{A6: FTZ/DAZ Apagado por Defecto --- IEEE-754 Conformance}

\subsection{El Problema}

V762 habilitaba FTZ/DAZ (Flush-To-Zero / Denormals-Are-Zero) incondicionalmente
en el \texttt{FtzDazGuard}. Esto rompe el estándar IEEE-754 y puede producir
resultados diferentes entre plataformas (x86 vs ARM64 vs código sin FTZ).

\subsection{La Solución V764}

\begin{verbatim}
/* A6: FTZ/DAZ apagado por defecto. IEEE-754 conforme salvo opt-in explícito. */
#ifndef POLYDIM_ENABLE_FTZ
  #define POLYDIM_ENABLE_FTZ 0
#endif

inline void set_fp_mode() {
#if POLYDIM_ENABLE_FTZ && defined(POLYDIM_X86)
    _MM_SET_FLUSH_ZERO_MODE(_MM_FLUSH_ZERO_ON);
    _MM_SET_DENORMALS_ZERO_MODE(_MM_DENORMALS_ZERO_ON);
#endif
    // Por defecto: IEEE-754 estricto. Sin flush de subnormales.
}
\end{verbatim}

El cambio tiene una consecuencia: el nuevo parámetro \texttt{reject\_subnormal = 0}
en las tolerancias por defecto significa que los subnormales son tolerados (no rechazados).
Un vector unitario disperso puede tener componentes subnormales legítimas que representan
contribuciones infinitesimales a la norma.

\section{A9: Umbral de Pivote Relativo a $\|{M}\|_\infty$}

En la retracción Cayley-SMW, la inversión de la matriz $\mathbf{M} \in \mathbb{R}^{2K \times 2K}$
falla si algún pivote es numéricamente cero. V762 usaba un umbral absoluto
\texttt{tol.pivot\_abs = 1e-14}.

El problema: para matrices con $\|{M}\|_\infty = 10^{-8}$ (caso de paso de aprendizaje muy pequeño),
todos los pivotes son pequeños pero la matriz está bien condicionada. Un umbral absoluto
rechazaría un caso perfectamente válido.

V764 usa umbral relativo:
\begin{equation}
\text{pivot\_threshold} = \text{tol.pivot\_rel} \times \|{M}\|_\infty = 8\varepsilon_{\text{mach}} \times \|{M}\|_\infty
\end{equation}

\section{A10: num\_threads Explícito}

V762 usaba \texttt{\#pragma omp parallel reduction} sin especificar \texttt{num\_threads},
pero el arreglo de acumuladores se indexaba por \texttt{omp\_get\_thread\_num()}. Si el entorno
variaba el número de threads entre el \texttt{parallel} externo y el interno, el índice podría
exceder el tamaño del arreglo.

V764 captura el número de threads una vez:
\begin{verbatim}
const int nthreads = clamp_threads(omp_get_max_threads());
std::vector<Neumaier> a_yy(nthreads), a_yu(nthreads), ...;
#pragma omp parallel num_threads(nthreads) // EXPLÍCITO en cada región
\end{verbatim}

\section{A11: Autodiagnóstico Anti-\texttt{-ffast-math}}

El flag de compilación \texttt{-ffast-math} habilita reasociación de operaciones de punto
flotante, lo que destruye el acumulador de Neumaier. La resta $(s - t)$ en el algoritmo
se vuelve algebraicamente cero bajo reasociación y el compilador la elimina.

V764 incluye una función de autodiagnóstico que falla si el compilador reasocio:

\begin{verbatim}
extern "C" POLYDIM_EXPORT int32_t POLYDIM_CALL polydim_selftest_compensation(void) {
    // Construir un caso donde Neumaier difiere de suma naïve en > 0.5 ULP
    // Si el compilador reasocio, ambas sumas darán lo mismo: FALLA
    const double big = 1e15, small = 1.0;
    Neumaier acc;
    acc.add(big); acc.add(-big); acc.add(small);
    const double naive = (big + (-big)) + small; // = 1.0 (si no hay reasoc)
    const double neumaier = acc.total();
    // Con reasociación: acc.total() = 0 (o NaN)
    // Sin reasociación: acc.total() = 1.0
    if (std::abs(neumaier - small) > 0.5 * std::numeric_limits<double>::epsilon())
        return POLYDIM_ERR_COMPENSATION_BROKEN; // -12 nuevo
    return POLYDIM_SUCCESS;
}

extern "C" POLYDIM_EXPORT int32_t POLYDIM_CALL polydim_selftest_all(void) {
    int32_t rc;
    if ((rc = polydim_selftest_compensation()) != POLYDIM_SUCCESS) return rc;
    // ... más tests ...
    return POLYDIM_SUCCESS;
}
\end{verbatim}

\textbf{[CERTIFICADO]} 
\textbf{Consecuencia para el Dart FFI (A12):} La función \texttt{polydim\_selftest\_all()}
es llamada automáticamente en \texttt{Polydim.open()}, garantizando que cualquier
binario compilado con \texttt{-ffast-math} sea detectado y rechazado \textit{antes}
de que ningún cálculo incorrecto salga del proceso.


\section{A12: El Bug más Costoso --- El Exit Code 0 Falso del Dart FFI}

\subsection{El Error que Invalida la Certificación de V761}

El Dart FFI de V761 tenía el siguiente main():

\begin{verbatim}
// V761 polydim_ffi_v761.dart - main() ROTO
void main() {
  final lib = DynamicLibrary.open('bin/polydim_kernel.so');
  final rodrigues = lib.lookupFunction<_RodriguesNative, _RodriguesDart>('polydim_rodrigues');
  // ^ Solo verifica que dlopen y dlsym funcionan
  // NUNCA llama a rodrigues(). 
  // NUNCA mide latencia.
  // El comentario "46 ms" era un literal en el código, no una medición.
  print('Exit 0'); // ← Esto es lo único que "certificaba"
}
\end{verbatim}

Adicionalmente:
\begin{enumerate}
\item El nombre de biblioteca era \texttt{bin/polydim\_kernel.so}: incorrecto en Linux/Android
(debe ser \texttt{libpolydim.so}) y sin soporte para macOS (\texttt{libpolydim.dylib}).
\item No se liberaba memoria: cada llamada filtraba $4 \times D \times 8$ bytes.
\item Sin \texttt{isLeaf}: sobrecarga de $235\,\text{ns}$ en lugar de $28\,\text{ns}$ por llamada.
\end{enumerate}

\subsection{La Corrección V764}

\begin{verbatim}C}, caption={Dart FFI V764 --- Medición real}, label=lst:dart_v764]
// V764 polydim_ffi_v764.dart - mide DE VERDAD
static String _defaultLibraryName() {
  if (Platform.isWindows) return 'polydim.dll';
  if (Platform.isMacOS) return 'libpolydim.dylib'; // A12.3: faltaba
  return 'libpolydim.so';  // Linux/Android: prefijo `lib`
}

// Descarga exhaustiva de candidatos de path
final candidates = [name, './$name', '$cwd/$name', '$cwd/build/$name'];

// Autodiagnóstico obligatorio al abrir:
static Polydim open({String? path, bool runSelftest = true}) {
  final p = Polydim._(lib);
  if (runSelftest) {
    final rc = p._selftestAll(); // A11: detecta -ffast-math
    if (rc != polydimSuccess) throw PolydimException('selftest_all', rc);
  }
  return p;
}

// Medición real con 9 repeticiones + mediana:
for (var i = 0; i < 9; i++) {
  final sw = Stopwatch()..start();
  rc = poly._rodrigues(pyn, pu, pv, po, 0.7, d, nullptr, rep);
  sw.stop();
  times.add(sw.elapsedMicroseconds / 1000.0);
}
times.sort();
print('mediana=${times[times.length ~/ 2].toStringAsFixed(2)} ms');
// ^ Número medido, no un literal. Medido: ~3.4 ms en D=1e6 con 2 hilos.
\end{verbatim}

\section{La Nueva API de Tolerancias y Reporte}

V764 introduce un sistema formal de tolerancias y reportes que permite al llamante
ver exactamente qué ocurrió dentro del kernel:

\begin{tabular}{llll}}
\hline
\textbf{Campo} & \textbf{V762} & \textbf{V764} \\
\hline
\texttt{basis\_ortho} & No existía & $64\varepsilon_{\text{mach}}$, independiente de $D$ \\
\texttt{point\_norm} & No existía & $64\varepsilon_{\text{mach}}$ \\
\texttt{gram\_ortho} & \texttt{tol = 2D*eps} (erróneo) & $64\varepsilon_{\text{mach}}\sqrt{D}$ (Stiefel) \\
\texttt{pivot\_rel} & Absoluto \texttt{1e-14} & Relativo $8\varepsilon_{\text{mach}} \cdot \|{M}\|_\infty$ \\
\texttt{reject\_subnormal} & Siempre rechazaba & Configurable (off por defecto) \\
\hline
\texttt{point\_norm\_err} & No existía & Error medido de $\|{\mathbf{y}\|}_2 - 1$ \\
\texttt{basis\_uu\_err} & No existía & $|\|{\mathbf{u}\|}^2 - 1|$ \\
\texttt{basis\_uv\_err} & No existía & $|\langle\mathbf{u},\mathbf{v}\rangle|$ \\
\texttt{out\_norm\_err} & No existía & Error medido de $\|{\mathbf{y}\|_{\text{out}}}_2 - 1$ \\
\texttt{threads\_used} & No existía & Número de hilos OpenMP usados \\
\hline


\end{tabular}

\section{El POLYDIM\_MAX\_K Corregido: De 1024 a 512}

V762 declaraba soporte hasta $K = 1024$ en la retracción Cayley-SMW. El hallazgo A1 reveló:

\begin{equation}
\text{Memoria para Gram } 2K \times 2K = (2 \times 1024)^2 \times 8 \text{ bytes} = 33.55\,\text{MB}
\end{equation}
\begin{equation}
\text{FLOPs para invertir } 2K \times 2K = \mathcal{O}((2K)^3) = \mathcal{O}(8 \times 10^9) \approx 8\,\text{GFLOP}
\end{equation}

$8\,\text{GFLOP}$ secuenciales para la inversión hacen que el tiempo del Cayley-SMW
sea dominado por la fase de pequeña matriz para $K = 1024$, negando el ahorro de escalar
en $D$. El máximo sostenible es $K = 512$:

\begin{equation}
\text{POLYDIM\_MAX\_K} = 512 \Rightarrow \text{Gram} = 1024 \times 1024 \times 8\,\text{B} = 8\,\text{MB}, \quad \text{inv} \approx 1\,\text{GFLOP}
\end{equation}

\section{Nuevos Códigos de Error en V764}

\begin{tabular}{llll}
\hline
\textbf{Código} & \textbf{Nombre} & \textbf{Causa} \\
\hline
$-8$ & \texttt{INVALID\_SCALAR} & \texttt{theta} o \texttt{tau} no finito (A3) \\
$-9$ & \texttt{BASIS\_NOT\_ORTHONORMAL} & Base $(\mathbf{u}, \mathbf{v})$ no ortonormal (A2) \\
$-10$ & \texttt{POINT\_OFF\_MANIFOLD} & $\mathbf{y} \notin S^{D-1}$ (A4) \\
$-11$ & \texttt{ALIASED\_BUFFERS} & \texttt{y\_out} solapa con \texttt{u} o \texttt{v} (A2) \\
$-12$ & \texttt{COMPENSATION\_BROKEN} & Compilador reasocio Neumaier (A11) \\
\hline


\end{tabular}

\textbf{[CERTIFICADO]} 
\textbf{Estado V764 (2026-09-20):} Los 12 hallazgos A1--A12 han sido corregidos
en el kernel. Los benchmarks base permanecen: Drift = $4.44 \times 10^{-16}$,
PMTP MaxBitDiff = 0. La API es ahora \textbf{más estricta} --- vectores no normalizados
son rechazados explícitamente en lugar de producir resultados silenciosamente incorrectos.

